% GENERATED FILE. Edit canonical lessons, then run npm run generate:books.
% canonical-source-hash: fnv1a64:967bbcf09fb25503
% canonical-lessons: BN-C358-mil, BN-C358-parthokyo, BN-C358-jogsutro, BN-C358-sadharon, BN-C358-guliye-phela

\chapter{Alike and Different}
\label{ch:bn-alike-and-different}
\hlchaptermodality{358}

\begin{chapteropening}
\textbf{By the end of this chapter:} \emph{I can say how two things are alike and how they differ.}

Complete the last lesson of chapter 358: I can say how two things are alike and how they differ.
\end{chapteropening}

\section[\texorpdfstring{\bn{মিল} (mil)}{মিল (mil)}]{\bn{মিল} (mil) — a likeness; a match}

\label{lesson:BN-C358-mil}



\begin{quote}
Before the new one: say the Bengali for a conclusion of a text, then the Bengali for a notion.
\end{quote}

\subsection*{You'll want to know: \bn{মিল}}
\textbf{\bn{মিল}} — \emph{mil} — \textquotedblleft{}a likeness; a match\textquotedblright{}.

\begin{quote}
\bn{দুই} \bn{ভাইয়ের} \bn{চেহারায়} \bn{অনেক} \bn{মিল} \bn{আছে।}
\end{quote}

\begin{quote}
\emph{dui bhāiyer chehārāy ônek mil āchhe.}
\end{quote}

\begin{quote}
The two brothers look very alike.
\end{quote}

\begin{cousinweb}
\textbf{\bn{মিল}} (\emph{mil}) is of the same family as \bn{মেলানো}, from Sanskrit \emph{milana}, meeting (CDIAL 10134 lists Bengali milan): the place where two things meet and agree.
\end{cousinweb}

\begin{grammarlens}[title={What you've built}]
One more word for giving reasons and opinions.
\end{grammarlens}

\subsection*{Your turn}
\begin{itemize}
\raggedright
  \item \emph{Say it:} \emph{mil}
  \item \emph{Say it:} \emph{mil}, once more
  \item \emph{Say it:} \emph{dui bhāiyer chehārāy ônek mil āchhe.} — The two brothers look very alike.
  \item \emph{Recall:} say the Bengali for a conclusion of a text, then the Bengali for a notion, then say \emph{mil} again
\end{itemize}

\begin{tcolorbox}[breakable,colback=teal!4,colframe=teal!35!black,title={Before you move on}]
What does \bn{মিল} mean? (\textquotedblleft{}A likeness; a match\textquotedblright{}.) Say it once more.
\end{tcolorbox}
\section[\texorpdfstring{\bn{পার্থক্য} (pārthôkyô)}{পার্থক্য (pārthôkyô)}]{\bn{পার্থক্য} (pārthôkyô) — a difference}

\label{lesson:BN-C358-parthokyo}



\begin{quote}
Before the new one: say the Bengali for a notion, then the Bengali for a likeness.
\end{quote}

\subsection*{You'll want to know: \bn{পার্থক্য}}
\textbf{\bn{পার্থক্য}} — \emph{pārthôkyô} — \textquotedblleft{}a difference\textquotedblright{}.

\begin{quote}
\bn{এই} \bn{দুই} \bn{শব্দের} \bn{পার্থক্য} \bn{কী}?
\end{quote}

\begin{quote}
\emph{ei dui shôbder pārthôkyô kī?}
\end{quote}

\begin{quote}
What is the difference between these two words?
\end{quote}

\begin{cousinweb}
\textbf{\bn{পার্থক্য}} (\emph{pārthôkyô}) is Sanskrit \emph{pārthakya}, difference (Monier-Williams), from \emph{pṛthak} 'apart'.
\end{cousinweb}

\begin{grammarlens}[title={What you've built}]
One more word for giving reasons and opinions.
\end{grammarlens}

\subsection*{Your turn}
\begin{itemize}
\raggedright
  \item \emph{Say it:} \emph{pārthôkyô}
  \item \emph{Say it:} \emph{pārthôkyô}, once more
  \item \emph{Say it:} \emph{ei dui shôbder pārthôkyô kī?} — What is the difference between these two words?
  \item \emph{Recall:} say the Bengali for a notion, then the Bengali for a likeness, then say \emph{pārthôkyô} again
\end{itemize}

\begin{tcolorbox}[breakable,colback=teal!4,colframe=teal!35!black,title={Before you move on}]
What does \bn{পার্থক্য} mean? (\textquotedblleft{}A difference\textquotedblright{}.) Say it once more.
\end{tcolorbox}
\section[\texorpdfstring{\bn{যোগসূত্র} (jogsūtrô)}{যোগসূত্র (jogsūtrô)}]{\bn{যোগসূত্র} (jogsūtrô) — a link; a connection}

\label{lesson:BN-C358-jogsutro}



\begin{quote}
Before the new one: say the Bengali for a likeness, then the Bengali for a difference.
\end{quote}

\subsection*{You'll want to know: \bn{যোগসূত্র}}
\textbf{\bn{যোগসূত্র}} — \emph{jogsūtrô} — \textquotedblleft{}a link; a connection\textquotedblright{}.

\begin{quote}
\bn{এই} \bn{দুই} \bn{ঘটনার} \bn{কোনো} \bn{যোগসূত্র} \bn{আছে}?
\end{quote}

\begin{quote}
\emph{ei dui ghôṭônār kono jogsūtrô āchhe?}
\end{quote}

\begin{quote}
Is there any link between these two events?
\end{quote}

\begin{cousinweb}
\textbf{\bn{যোগসূত্র}} (\emph{jogsūtrô}) is Sanskrit \emph{yoga}, joining, with \emph{sūtra}, a thread (Monier-Williams): the thread that ties two things together.
\end{cousinweb}

\begin{grammarlens}[title={What you've built}]
One more word for giving reasons and opinions.
\end{grammarlens}

\subsection*{Your turn}
\begin{itemize}
\raggedright
  \item \emph{Say it:} \emph{jogsūtrô}
  \item \emph{Say it:} \emph{jogsūtrô}, once more
  \item \emph{Say it:} \emph{ei dui ghôṭônār kono jogsūtrô āchhe?} — Is there any link between these two events?
  \item \emph{Recall:} say the Bengali for a likeness, then the Bengali for a difference, then say \emph{jogsūtrô} again
\end{itemize}

\begin{tcolorbox}[breakable,colback=teal!4,colframe=teal!35!black,title={Before you move on}]
What does \bn{যোগসূত্র} mean? (\textquotedblleft{}A link; a connection\textquotedblright{}.) Say it once more.
\end{tcolorbox}
\section[\texorpdfstring{\bn{সাধারণ} (sādhārôn)}{সাধারণ (sādhārôn)}]{\bn{সাধারণ} (sādhārôn) — ordinary; common}

\label{lesson:BN-C358-sadharon}



\begin{quote}
Before the new one: say the Bengali for a difference, then the Bengali for a link.
\end{quote}

\subsection*{You'll want to know: \bn{সাধারণ}}
\textbf{\bn{সাধারণ}} — \emph{sādhārôn} — \textquotedblleft{}ordinary; common\textquotedblright{}.

\begin{quote}
\bn{আমি} \bn{একজন} \bn{সাধারণ} \bn{মানুষ।}
\end{quote}

\begin{quote}
\emph{āmi ekjôn sādhārôn mānush.}
\end{quote}

\begin{quote}
I'm an ordinary person.
\end{quote}

\begin{cousinweb}
\textbf{\bn{সাধারণ}} (\emph{sādhārôn}) is Sanskrit \emph{sādhāraṇa}, common to all (Monier-Williams). In Bengali it is both 'general' and 'ordinary, nothing special'.
\end{cousinweb}

\begin{grammarlens}[title={What you've built}]
One more word for giving reasons and opinions.
\end{grammarlens}

\subsection*{Your turn}
\begin{itemize}
\raggedright
  \item \emph{Say it:} \emph{sādhārôn}
  \item \emph{Say it:} \emph{sādhārôn}, once more
  \item \emph{Say it:} \emph{āmi ekjôn sādhārôn mānush.} — I'm an ordinary person.
  \item \emph{Recall:} say the Bengali for a difference, then the Bengali for a link, then say \emph{sādhārôn} again
\end{itemize}

\begin{tcolorbox}[breakable,colback=teal!4,colframe=teal!35!black,title={Before you move on}]
What does \bn{সাধারণ} mean? (\textquotedblleft{}Ordinary; common\textquotedblright{}.) Say it once more.
\end{tcolorbox}
\section[\texorpdfstring{\bn{গুলিয়ে} \bn{ফেলা} (guliye phelā)}{গুলিয়ে ফেলা (guliye phelā)}]{\bn{গুলিয়ে} \bn{ফেলা} (guliye phelā) — to mix up; to confuse (two things)}

\label{lesson:BN-C358-guliye-phela}



\begin{quote}
Before the new one: say the Bengali for a link, then the Bengali for ordinary.
\end{quote}

\subsection*{You'll want to know: \bn{গুলিয়ে} \bn{ফেলা}}
\textbf{\bn{গুলিয়ে} \bn{ফেলা}} — \emph{guliye phelā} — \textquotedblleft{}to mix up; to confuse (two things)\textquotedblright{}.

\begin{quote}
\bn{আমি} \bn{প্রায়ই} \bn{এই} \bn{দুই} \bn{শব্দ} \bn{গুলিয়ে} \bn{ফেলি।}
\end{quote}

\begin{quote}
\emph{āmi prāyi ei dui shôbdo guliye pheli.}
\end{quote}

\begin{quote}
I often mix up these two words.
\end{quote}

Past (he/she …): \bn{তিনি} \bn{নামগুলো} \bn{গুলিয়ে} \bn{ফেললেন}, \emph{tini nāmgulo guliye phellen}.

\begin{cousinweb}
\textbf{\bn{গুলিয়ে} \bn{ফেলা}} (\emph{guliye phelā}) is the -\bn{এ} form of \bn{গোলা}, to dissolve or stir into water, with \bn{ফেলা} 'to throw' (CDIAL 8651 lists Bengali phelā), which marks a thing done and finished: two things stirred together until you cannot tell them apart. The source of \bn{গোলা} is not settled; CDIAL lists Bengali gholā 'to stir, mix' (4526) and galā 'to melt' (4074), and neither is certain.
\end{cousinweb}

\begin{grammarlens}[title={What you've built}]
One more word for giving reasons and opinions.
\end{grammarlens}

\subsection*{Your turn}
\begin{itemize}
\raggedright
  \item \emph{Say it:} \emph{guliye phelā}
  \item \emph{Say it:} \emph{guliye phelā}, once more
  \item \emph{Say it:} \emph{āmi prāyi ei dui shôbdo guliye pheli.} — I often mix up these two words.
  \item \emph{Recall:} say the Bengali for a link, then the Bengali for ordinary, then say \emph{guliye phelā} again
\end{itemize}

\begin{tcolorbox}[breakable,colback=teal!4,colframe=teal!35!black,title={Before you move on}]
What does \bn{গুলিয়ে} \bn{ফেলা} mean? (\textquotedblleft{}To mix up; to confuse (two things)\textquotedblright{}.) Say it once more.
\end{tcolorbox}
